\chapter{Практика 5. OFDM модуляция}
\label{ch:chap1}

\section{Задание}

Реализовать формирование OFDM символа, включающего в себя информационные символы и опорные сигналы.

[........................................................................................................................]

\section{Описание работы модулятора}

Цель применения технологии OFDM состоит в борьбе с многолучевостью. В реализуемой имитационной модели задача блока OFDM модулятора заключается в формировании OFDM символа, включающего информационные символы и опорный сигнал. Опорный сигнал потребуется для оценки АЧХ канала передачи при эквалайзировании сигнала в приемнике. В информационных символах содержится передаваемое текстовое сообщение.

Соответственно, на вход модулятора должны поступать информационные QPSK символы (вектор комплексных отсчетов), а на выходе — IQ отсчеты OFDM символа (вектор комплексных отсчетов).

\begin{figure}[H]
\centering
\includegraphics[width=0.7\textwidth]{ofdm.png}
\caption{Схема OFDM модулятора}
\end{figure}

Рассмотрим подробнее блок модели. OFDM модулятор включает в себя несколько последовательных процедур.

\subsubsection{1. Добавление опорного сигнала}

Эта процедура включает в себя последовательное выполнение следующих операций:

\textbf{а)} Расчет количества опорных поднесущих в зависимости от шага размещения, который обозначен здесь переменной $\Delta{RS}$.
В модели должен быть реализован алгоритм расположения опорных под-
несущих между поднесущими с информационными данными. При этом, период (или шаг) размещения
опорных поднесущих $\Delta{RS}$ должен быть контролируемым
пользователем, т.е. пользователь должен иметь возможность его изменять в программе. 

\begin{figure}[H]
\centering
\includegraphics[width=0.7\textwidth]{podnesysh.png}
\caption{Алгоритм расположения поднесущих}
\end{figure}

Общее количество опорных поднесущих в OFDM 
символе обозначим переменной $N_{RS}$. Переменная $N_{RS}$ напрямую зависит
от $\Delta{RS}$, потому что чем меньше шаг размещения, тем боле опорных поднесущих будет в OFDM символе и, наоборот. Для расчета $N_{RS}$ может быть
использована следующая формула:

\[
N_{RS} = \left\lfloor \frac{N_{QPSK}}{\Delta_{RS}} \right\rfloor
\]

где $N_{QPSK}$ — количество QPSK символов на выходе модулятора. а скобки [∙] означают округление до ближайшего целого числа в меньшую сторону.

\textbf{б)} Формирование опорного сигнала. В качестве опорного сигнала могут использоваться любые комплексные сигнальные 
отсчеты количеством $N_{RS}$ вида 
$I + jQ$ мощность которых нормирована на 1 Вт, т.е.
\[
\sqrt{I^2 + Q^2} = 1
\]
Например, для опорного сигнала можно сформировать
вектор длиной $N_{RS}$, состоящий из одинаковых отсчетов

\[
0.707 + j0.707
\]

\textbf{в)} Расчет индексов опорных поднесущих. 
Индексы поднесущих для $RS$ должны быть рассчитаны с учетом шага $\Delta{RS}$, указанного
пользователем. Например, если $\Delta{RS}$ = 6, то опорные поднесущие будут
располагаться через 6 поднесущих, а в этом случае индексы опорных поднесущих будут: 1, 7, 13,... Напоминаю, что общее количество индексов
(как и опорных поднесущих) должно быть равно $N_{RS}$.

\textbf{г)} Размещение опорного сигнала, 
сформированного на этапе (б) по опорным поднесущим с индексами, рассчитанными на этапе (в).

\subsubsection{2. Добавление информационного сигнала}

Чтобы реализовать эту
процедуру, нужно рассчитать индексы поднесущих для информационного сигнала.

\subsubsection{3. Добавление нулевого защитного интервала слева и справа от
информационной части}

Задача защитного интервала – борьба с внеполосным излучением. Требуется
добавить в спектр OFDM символа нулевые поднесущие слева и справа от
активных поднесущих (которые используются для передачи информаци-
онного и опорного сигнала). 

\begin{figure}[H]
\centering
\includegraphics[width=0.7\textwidth]{null.png}
\caption{Добавление нулевого защитного интервала}
\end{figure}

Количество нулевых поднесущих $N_{Z}$ с каж-
дой стороны определяется параметром $C$ по формуле:

\[
N_Z = C \cdot (N_{RS} + N_{QPSK})
\]

Параметр $C$ отличается для разных вариантов задания. Например, если
$C = 1/4$, $N_{RS} = 10$, $N_{QPSK} = 90$ то количество нулей слева и справа будет $N_{Z} = 25$.

\subsubsection{4. Обратное дискретное преобразование Фурье} 
ОДПФ над всеми поднесущими для формирования OFDM символа во временном домене. Размерность ОДПФ должна быть равна общему количеству всех поднесущих
(включая нулевые, опорные и информационные).

\subsubsection{5. Добавление циклического префикса}

Задача циклического префикса – борьба с межсимвольной интерференцией. Циклический 
префикс является копией отрезка конца OFDM
символа длительностью $T_{CP}$, помещенной в начало символа. При этом
не нарушается ортогональность за счет свойств ДПФ.

\begin{figure}[H]
\centering
\includegraphics[width=0.5\textwidth]{prefix.png}
\caption{Добавление циклического префикса}
\end{figure}


Сигнал на выходе блока OFDM модулятора должен являться вектором с комплексными отсчетами длиной:
\[
T_{CP} + N_{RS} + N_{QPSK} + 2 \cdot N_{Z}
\]

Процедура формирования OFDM символа визуализирована на рисунке ниже. Здесь канальный мультиплексор отвечает за расстановку сигналов (нулей, опорных сигналов и сигналов с данными) по поднесущим в
частотном домене.

\begin{figure}[H]
\centering
\includegraphics[width=0.8\textwidth]{ofdm_symbol.png}
\caption{Процедура формирования OFDM символа}
\end{figure}

\subsection{Описание входных/выходных переменных}

\textbf{Входные переменные:}
\begin{itemize}
\item вектор с комплексными символами
\end{itemize}

\textbf{Выходные переменные:}
\begin{itemize}
\item вектор с комплексными символами
\end{itemize}

\textbf{Параметры:}
\begin{itemize}
\item параметры OFDM символа
\item $\Delta{RS}$ - шаг опорных поднесущих
\item $T$ - размер циклич. префикса
\end{itemize}

\endinput